% Section 3.8, translated from sv/chapters/03/exercises.tex.

\section{Exercises}
\label{c3:sec:uppgifter}

\begin{aside}
\textbf{How to work with the exercises.} Parts~1 and~2 are done during the lesson: first on your
own, a few minutes per exercise, then together. Part~3 is extra exercises for further practice
after the lesson. Most of them can be done in your head or with pencil and paper.

\facitnotis{L03}
\end{aside}

% ------------------------------------------------------------------------------------------
\uppgiftsdel{Part 1}{Review exercises}

\uppgift{1.1}{Simplification}
Simplify the following expressions.
\begin{delar}(3)
  \task $3(2R + 4) - 2(R - 1)$
  \task $(U + 3)^2 - 9$
  \task $\dfrac{6I^2 + 4I}{2I}$, assume $I \neq 0$
\end{delar}

\uppgift{1.2}{Factorisation}
Factorise.
\begin{delar}(3)
  \task $R^2 - 25$
  \task $3U^2 + 6U$
  \task $I^2 - 8I + 16$
\end{delar}

% ------------------------------------------------------------------------------------------
\needspace{16\baselineskip}
\uppgiftsdel{Part 2}{New material}

\uppgift{2.1}{Ohm's law: find the resistance}
The voltage across a resistor is $U = 15\,\text{V}$ and the current is $I = 0.3\,\text{A}$.

\bild[0.31\linewidth]{lectures/L03/appendix/images/2.1_circuit.png}

Ohm's law states that $U = R \cdot I$.
\begin{delar}(2)
  \task Set up the equation and solve for $R$.
  \task Check your answer.
\end{delar}

\uppgift{2.2}{Series circuit: an equation with brackets}
Two resistors are connected in series. The circuit is supplied with $U = 24\,\text{V}$ and the
current is $I = 0.4\,\text{A}$. One resistor is $R_2 = 40\,\Omega$; the other is unknown.

\bild[0.41\linewidth]{lectures/L03/appendix/images/2.2_circuit.png}

Ohm's law for the whole circuit gives
\[
  U = (R_1 + R_2) \cdot I.
\]
\begin{parts}
  \item Substitute the values and solve the equation for $R_1$ \textbf{without} expanding the
    brackets.
  \item Solve the equation again, but expand the brackets first. Check that you get the same
    answer.
  \item Check the solution by calculating the current $I = \dfrac{U}{R_1 + R_2}$.
\end{parts}

\uppgift{2.3}{Voltage division}
In a voltage divider, the output voltage is given by
\[
  U_{\text{out}} = \frac{R_2}{R_1 + R_2} \cdot U_{\text{in}}.
\]

\bild[0.44\linewidth]{lectures/L03/appendix/images/2.3_circuit_en.png}

Assume that $U_{\text{in}} = 9\,\text{V}$, $R_2 = 3\,\Omega$ and $U_{\text{out}} = 3\,\text{V}$.
\begin{parts}
  \item Set up the equation for $U_{\text{out}}$ and solve for $R_1$.
  \item Check your answer.
\end{parts}

\uppgift{2.4}{Time constant of an RC circuit}
The time constant $\tau$ (tau) of an RC circuit is
\[
  \tau = R \cdot C.
\]

\bild[0.39\linewidth]{lectures/L03/appendix/images/2.4_circuit.png}

The capacitor's capacitance is $C = 100\,\upmu\text{F} = 100 \times 10^{-6}\,\text{F}$ and the
time constant is $\tau = 0.05\,\text{s}$. Calculate the resistance $R$.

\needspace{16\baselineskip}
\uppgift{2.5}{Safe operating area}
The power dissipation of a transistor is calculated as $P = U_{\text{CE}} \cdot I_{\text{C}}$.

\bild[0.4\linewidth]{lectures/L03/appendix/images/2.5_circuit_en.png}

The transistor can take at most $P_{\max} = 500\,\text{mW}$, and the collector current is
$I_{\text{C}} = 50\,\text{mA}$.
\begin{parts}
  \item Set up an inequality to find the largest permitted voltage drop $U_{\text{CE}}$.
  \item Solve the inequality and give the largest permitted $U_{\text{CE}}$ in volts.
\end{parts}

\uppgift{2.6}{Loaded battery: an inequality where the sign reverses}
A battery has the open-circuit voltage $E = 12\,\text{V}$ and the internal resistance
$R_{\text{i}} = 0.8\,\Omega$. When the battery is loaded with the current $I$, the terminal
voltage is
\[
  U = E - R_{\text{i}} \cdot I.
\]

\bild[0.43\linewidth]{lectures/L03/appendix/images/2.6_circuit.png}

\begin{parts}
  \item Calculate the terminal voltage $U$ when $I = 2\,\text{A}$.
  \item The connected device needs at least $9\,\text{V}$ to work. Set up the inequality
    $U \geq 9$ and solve it for $I$.
    \begin{aside}
    \note[Tip] Think about what happens to the inequality sign when you divide by a negative
    number.
    \end{aside}
  \item Another battery has $E = 13\,\text{V}$ and $R_{\text{i}} = 1.0\,\Omega$. At what current
    do the two batteries give the same terminal voltage? Set up an equation with the unknown on
    both sides and solve it.
  \item Which of the batteries gives the higher terminal voltage at $I = 8\,\text{A}$?
\end{parts}

\uppgift{2.7}{Absolute value in signal analysis}
A signal can deviate positively or negatively from a reference voltage. The deviation $\delta$
satisfies $\abs{\delta} \leq 0.5\,\text{V}$.
\begin{parts}
  \item Rewrite the absolute-value inequality as a double inequality without the absolute value.
  \item If the reference voltage is $5\,\text{V}$, give the permitted range of the actual
    voltage.
\end{parts}

\uppgift{2.8}{Current regulation: an absolute-value equation}
A current regulator compares the measured current $I$ (in amperes) with its setpoint and gives
the error voltage
\[
  u_{\text{e}} = 0.4I - 2.
\]
\begin{parts}
  \item At what current is the error voltage zero, that is, what is the setpoint?
  \item The regulator raises an alarm when the error voltage has the magnitude $1.2\,\text{V}$.
    Solve the equation $\abs{0.4I - 2} = 1.2$ as two cases and give the two currents at which the
    alarm is triggered.
  \item Which range of currents satisfies $\abs{0.4I - 2} < 1.2$, that is, when does the
    regulator \emph{not} raise an alarm?
\end{parts}

% ------------------------------------------------------------------------------------------
\uppgiftsdel{Part 3}{Extra exercises}
The exercises below are extra practice, best done on your own after the lesson.

\uppgift{3.1}{Simple linear equations}
Solve the equations.
\begin{delar}(3)
  \task $x + 7 = 12$
  \task $5x = 45$
  \task $3x - 4 = 11$
  \task $\dfrac{x}{4} = 6$
  \task $8 - 2x = 0$
  \task $-3x = 18$
\end{delar}

\uppgift{3.2}{The unknown on both sides}
Solve the equations and check your answers.
\begin{delar}(3)
  \task $5x - 3 = 2x + 9$
  \task $7 - 2x = x + 1$
  \task $4(x - 1) = 2(x + 3)$
  \task $3(2x + 1) = 6x + 3$
  \task $2(x + 4) = 2x + 5$
\end{delar}

\begin{aside}
\note[Tip] Two of the equations behave differently from the others. Think about what that means.
\end{aside}

\uppgift{3.3}{Equations with fractions}
Solve the equations.
\begin{delar}(2)
  \task $\dfrac{x}{3} + \dfrac{x}{6} = 3$
  \task $\dfrac{2x - 1}{5} = 3$
  \task $\dfrac{x}{2} - \dfrac{x}{5} = 3$
  \task $\dfrac{12}{x} = 4$, assume $x \neq 0$
\end{delar}

\uppgift{3.4}{Proportions}
\begin{parts}
  \item Solve $\dfrac{x}{8} = \dfrac{3}{4}$.
  \item Solve $\dfrac{5}{x} = \dfrac{2}{6}$.
  \item In a voltage divider, $\dfrac{U_1}{U_2} = \dfrac{R_1}{R_2}$. Calculate $R_1$ when
    $U_1 = 4\,\text{V}$, $U_2 = 8\,\text{V}$ and $R_2 = 2\,\text{k}\Omega$.
  \item A resistor has $6\,\text{V}$ across it at the current $0.2\,\text{A}$. What voltage is
    needed for the current $0.5\,\text{A}$? Set up a proportion and solve it.
\end{parts}

\uppgift{3.5}{Solving a formula for a variable}
Solve for the variable given.
\begin{delar}(2)
  \task $U = RI$, solve for $I$.
  \task $P = RI^2$, solve for $R$.
  \task $R_{\text{TOT}} = R_1 + R_2$, solve for $R_2$.
  \task $\tau = RC$, solve for $C$.
  \task $U = E - R_{\text{i}}I$, solve for $R_{\text{i}}$.
  \task $\eta = \dfrac{P_{\text{out}}}{P_{\text{in}}}$, solve for $P_{\text{in}}$.
\end{delar}

\uppgift{3.6}{Inequalities}
Solve the inequalities and describe each solution as an interval on the number line.
\begin{delar}(3)
  \task $x + 3 < 8$
  \task $2x - 1 \geq 7$
  \task $-x > 4$
  \task $-2x + 6 \leq 0$
  \task $5 - 3x > 14$
\end{delar}

\uppgift{3.7}{Double inequalities}
\begin{parts}
  \item Solve $2 < x + 1 < 7$.
  \item Solve $-3 \leq 2x - 1 \leq 5$.
  \item A component is supplied with $U = 5\,\text{V}$ and can take at most
    $P_{\max} = 0.25\,\text{W}$. Set up and solve an inequality for the permitted current $I$.
  \item An input is read as a logic one when its voltage lies in the range
    $2.0\,\text{V} \leq u \leq 3.3\,\text{V}$. The input voltage is $u = 0.5x$ volts. Which
    values of $x$ give a logic one?
\end{parts}

\uppgift{3.8}{Absolute values in equations}
Solve the equations.
\begin{delar}(3)
  \task $\abs{x} = 6$
  \task $\abs{x - 3} = 5$
  \task $\abs{2x + 1} = 7$
  \task $\abs{x + 4} = 0$
  \task $\abs{3x - 2} = -5$
\end{delar}

\uppgift{3.9}{Absolute values in inequalities and tolerances}
\begin{parts}
  \item Write $\abs{x - 5} \leq 2$ without the absolute value and solve the inequality.
  \item Solve $\abs{2x - 6} < 4$.
  \item A supply voltage is to be $5\,\text{V}$ with a tolerance of $\pm 2\,\%$. Write the
    requirement as an absolute-value inequality and give the permitted voltage range.
  \item A resistor is marked $1\,\text{k}\Omega \pm 1\,\%$ and measures $1\,015\,\Omega$. Does
    the resistor meet the requirement ${\abs{R - 1000} \leq 10}$?
\end{parts}

\uppgift{3.10}{Temperature-dependent resistance}
The resistance of a resistor varies with temperature as
\[
  R = R_0(1 + \alpha \cdot \Delta T).
\]
Here $R_0 = 100\,\Omega$ is the resistance at $20\,{}^{\circ}\text{C}$,
$\alpha = 0.004\,{}^{\circ}\text{C}^{-1}$ and $\Delta T$ is the change in temperature from
$20\,{}^{\circ}\text{C}$.
\begin{parts}
  \item Calculate $R$ at $70\,{}^{\circ}\text{C}$.
  \item At what temperature is $R = 110\,\Omega$?
  \item The circuit works as long as $R \leq 130\,\Omega$. What temperature range is permitted?
  \item Solve the formula for $\Delta T$ in general.
\end{parts}
